\documentclass[12pt]{article}
\usepackage{amsmath}
\usepackage{graphicx}
\usepackage{hyperref}
\usepackage{geometry}
\geometry{a4paper, margin=1in}

\title{Congestion and Deadlock Prediction in Multi-Robot Systems}
\author{}
\date{\today}

\begin{document}
\maketitle
\tableofcontents
\newpage

\begin{abstract}
This report details the development of a predictive machine learning model for identifying potential deadlocks and congestion conflicts in multi-robot systems. By engineering scale-invariant spatial, density, and behavioral features, we train an XGBoost classifier to forecast conflicts up to five time steps in the future. The approach emphasizes generalization across different warehouse scenarios by relying on relative density metrics and local behavioral indicators rather than absolute coordinates.
\end{abstract}

\section{Introduction}

\subsection{Background}
In multi-agent pathfinding (MAPF) and automated warehouse environments, robots frequently encounter congestion that can lead to deadlocks or prolonged waiting times. Efficient traffic management requires identifying these bottlenecks proactively rather than reacting to them after they form.

\subsection{Problem Statement}
Given a snapshot of the warehouse state—including robot positions, states, and local density—the task is to predict whether a robot will experience a conflict (collision, stuck state, or deadlock) at a specific prediction horizon ($t+5$). This is framed as a binary classification problem (Normal Flow vs. Conflict Risk) using leading indicators without leaking the current conflict state.

\subsection{Project Objectives}
\begin{itemize}
    \item Perform robust exploratory data analysis on the raw congestion dataset.
    \item Design a preprocessing pipeline that extracts scale-invariant density features and behavioral early warnings.
    \item Prevent data leakage by carefully isolating spatial risk mapping to the training set and excluding current-state conflict indicators.
    \item Train and regularize an XGBoost classifier optimized for the highly imbalanced nature of the dataset.
    \item Evaluate the model's generalization capability across fundamentally different traffic scenarios.
\end{itemize}

\subsection{Report Structure}
Section 2 details the dataset schema, target variable, and exploratory analysis. Section 3 outlines the methodology, including the task formulation and experimental setup. Section 4 covers the comprehensive data preprocessing and feature engineering pipeline.

\section{Dataset \& Exploratory Data Analysis}

\subsection{Dataset Schema \& Overview}
The raw dataset (\texttt{congestion\_dataset.csv}) is generated by simulating multi-agent pathfinding (MAPF) operations within a grid-based automated warehouse. It captures a snapshot of every robot at each discrete time step ($t$). The recorded variables include spatial data (\texttt{x}, \texttt{y}), agent states (\texttt{state}, \texttt{held\_item}), local density counts (\texttt{local\_density\_3}, \texttt{local\_density\_5}), and conflict indicators (\texttt{is\_stuck}, \texttt{is\_deadlocked}, \texttt{collision\_vertex}).

\subsection{Target Variable Distribution}
The target variable is the future occurrence of a conflict ($t+5$). Because deadlocks and collisions are rare compared to the vast amount of normal robot flow, the target is highly imbalanced. This imbalance necessitates specialized metrics and dynamic weight scaling during model training.

\subsection{Numeric Feature Analysis}
The primary numeric features include spatial coordinates and absolute local densities. We observe that raw density counts perform poorly when applied to warehouses with different baseline populations. Consequently, the analysis highlights the need for relative, scale-invariant density transformations.

\subsection{Categorical \& Text Feature Analysis}
Categorical attributes include the robot's operational \texttt{state} (e.g., \texttt{AVAILABLE}, \texttt{IN\_PROGRESS}) and the \texttt{held\_item} field. The \texttt{held\_item} field is treated as a binary indicator denoting whether the robot is carrying payload, which influences routing priority and behavior.

\subsection{Missing Value Analysis}
While most features are fully populated, the predictive target shifting process generates missing labels at the very end of a robot's trajectory (where the future horizon $t+5$ is undefined). These trailing rows are dropped to maintain training integrity.

\subsection{Data Leakage Prevention}
A critical design decision is the exclusion of columns that describe the \textit{current} conflict state (e.g., \texttt{is\_stuck}, \texttt{is\_deadlocked}, \texttt{path\_length}). Including them would constitute target leakage, allowing the model to shortcut by recognizing an already-failed state rather than learning genuine early-warning behavioral precursors.

\section{Methodology}

\subsection{Task Formulation}
The task is formulated as a \textbf{binary classification} problem. We aim to classify each robot's current state into one of two risk categories:
\begin{itemize}
    \item \textbf{0 (Normal Flow):} No conflict predicted within the horizon.
    \item \textbf{1 (Conflict Risk):} The robot is predicted to be stuck, deadlocked, or in a collision at $t+5$.
\end{itemize}

\subsection{Test Model}

\subsubsection{Train-Test Split}
To rigorously evaluate predictive capability in a real-world scenario, we employ a strict 80/20 time-based split based on the simulation tick (\texttt{t}). Shuffling is intentionally avoided to preserve temporal causality and ensure that future target data never leaks into the training window.

\subsubsection{Model Selection}
For this predictive task, we utilize an XGBoost classifier. This model was chosen for its robustness on tabular data, its inherent ability to model complex, non-linear feature interactions (such as the relationship between local density peaks and stationary duration), and its computational efficiency. Dynamic \texttt{scale\_pos\_weight} is applied to handle the class imbalance.

\subsubsection{Evaluation Metrics}
Because of the heavy class imbalance, accuracy is an insufficient metric. Models are evaluated using:
\begin{itemize}
    \item \textbf{ROC-AUC:} To measure overall probabilistic ranking capability.
    \item \textbf{PR-AUC:} To evaluate the precision-recall trade-off on the minority risk class.
    \item \textbf{Optimal F1-Score:} The probability threshold is dynamically swept to find the point that maximizes the F1-score, providing a balanced, actionable decision boundary for early intervention.
\end{itemize}

\subsubsection{Cross-Dataset Generalization Results}
To validate that our feature engineering pipeline successfully captures generalized leading indicators rather than memorizing absolute densities, we performed a cross-dataset evaluation across three distinct environments: \texttt{few\_collisions} (baseline), \texttt{many\_collisions} (heavily congested), and \texttt{no\_collisions} (sparse, no conflicts).

\begin{table}[htbp]
\centering
\begin{tabular}{lccccccc}
\hline
\textbf{Experiment} & \textbf{Acc} & \textbf{ROC} & \textbf{PR} & \textbf{Thr} & \textbf{Prec} & \textbf{Rec} & \textbf{F1} \\
\hline
few $\rightarrow$ many & 58.0\% & 0.5865 & 0.6260 & 0.002 & 0.5708 & 0.9955 & 0.7256 \\
many $\rightarrow$ few & 68.3\% & 0.6900 & 0.6407 & 0.184 & 0.6404 & 0.8014 & 0.7119 \\
few+many $\rightarrow$ none & 40.2\% & 0.5736 & 0.4754 & 0.011 & 0.3936 & 0.9897 & 0.5632 \\
\hline
\end{tabular}
\caption{Cross-Dataset Evaluation Summary based on optimal F1-score thresholds.}
\end{table}

The results demonstrate the model's robustness to shifting traffic distributions:
\begin{itemize}
    \item \textbf{few $\rightarrow$ many collisions:} When moving to a heavily congested environment, the model automatically favors a very low threshold ($0.002$) to maximize recall ($99.55\%$). It successfully identifies almost all incoming deadlocks despite never seeing this level of congestion in training.
    \item \textbf{many $\rightarrow$ few collisions:} Conversely, when transitioning from a dense to a sparse environment, the scale-invariant features prevent the model from over-predicting. It maintains a balanced precision ($64.04\%$) and recall ($80.14\%$) with a strong F1 score of $0.7119$.
    \item \textbf{few+many $\rightarrow$ no collisions:} In a completely conflict-free dataset, the model acts conservatively. It flags high-risk behavioral states (leading to high recall of dangerous situations). The precision drops because these risky states did not ultimately culminate in actual collisions (false positives), reflecting a safety-first prediction behavior in novel environments.
\end{itemize}

\section{Data Preprocessing Pipeline}
To train a model capable of generalizing across different warehouse layouts and baseline congestion levels, the raw generated data undergoes a rigorous sequence of transformations. 

\subsection{Predictive Target Creation}
Instead of predicting current deadlocks, we construct a forward-looking predictive label (\texttt{target\_future\_conflict}). By applying a forward look-ahead window, if a robot is flagged as stuck, deadlocked, or colliding $N=5$ time steps in the future, the current time step is labeled as a positive risk.

\subsection{Categorical Encoding and Spatial Normalization}
Robot states are one-hot encoded, and \texttt{held\_item} is converted into a binary \texttt{is\_carrying} flag. To make the model agnostic to the absolute size of the warehouse grid, the raw $x$ and $y$ coordinates are min-max normalized into \texttt{x\_norm} and \texttt{y\_norm}.

\subsection{Scale-Invariant Density Features}
Raw density counts are transformed into scale-invariant features based on the robot's recent history:
\begin{itemize}
    \item \textbf{Density Z-Scores:} The z-score captures how unusually crowded the area is compared to the robot's immediate past (5-step window).
    \item \textbf{Density Peaks \& Acceleration:} We compute the rolling maximum (peak) to capture sudden spikes, and density acceleration (second difference) to indicate how rapidly congestion is building.
\end{itemize}

\subsection{Behavioral Early Warnings}
We engineer leading behavioral indicators that precede deadlocks without leaking the target:
\begin{itemize}
    \item \textbf{Position Stationarity:} We compute \texttt{stationary\_ratio}, the fraction of the last 5 steps spent without moving.
    \item \textbf{System-Wide Pressure:} \texttt{neighbour\_static\_count} counts how many \textit{other} robots in the system are currently frozen.
    \item \textbf{Density Spread:} The ratio of wider density (radius 5) to inner density (radius 3) measures whether robots are tightly clustered or broadly distributed.
\end{itemize}

\subsection{Spatial Risk Mapping}
Certain topological features (e.g., narrow corridors) are inherently more prone to congestion. A \textit{vertex risk factor} captures the historical risk of specific $(x, y)$ coordinates. Crucially, to prevent data leakage, this risk map is generated strictly from the historical distribution of conflicts in the training partition and deterministically mapped onto the test set.

\end{document}
